\section{Split Array -- Largest Sum}
\label{sec:bs-split-array}

\problemstatement{We are given an array $\textit{nums}$ of $n$ non-negative
integers and an integer $m$. We must split $\textit{nums}$ into $m$
\textbf{contiguous, non-empty} subarrays so as to \textbf{minimize the
largest sum} among the $m$ subarrays.}

\begin{patternbox}
This is, word for word, the Book Allocation Problem
(Section~\ref{sec:bs-book-allocation}) with ``books/pages/students''
relabeled as ``array/elements/subarrays.'' There is no new idea to extract
here beyond confirming, explicitly, that you recognize the relabeling --
which is itself the point of including this problem on the sheet
immediately after book allocation: to test whether the pattern has truly
been internalized, or only memorized under one specific cover story.
\end{patternbox}

\subsection*{Understanding the problem carefully}

Exactly as in Section~\ref{sec:bs-book-allocation}, define
$\textit{partsNeeded}(\textit{maxSum})$ as the minimum number of
contiguous parts needed if no part's sum may exceed $\textit{maxSum}$,
computed by the same greedy bucket-filling simulation. This is
non-increasing in $\textit{maxSum}$, giving the same monotonic structure.

\begin{ideabox}[Candidate Space and Predicate]
Candidate space: integer sums in $[\max(\textit{nums}),
\sum(\textit{nums})]$. Predicate:
$P(\textit{maxSum}) = (\textit{partsNeeded}(\textit{maxSum}) \le m)$. Want
the smallest $\textit{maxSum}$ where $P$ holds.
\end{ideabox}

\subsection*{Algorithm}

Identical to Algorithm~\ref{sec:bs-book-allocation}'s
\textsc{BookAllocation}, with \textit{pages} relabeled \textit{nums} and
$m$ meaning the number of parts directly.

\subsection*{Dry run}

\dryrun{Take $\textit{nums} = [7,2,5,10,8]$, $m = 2$.}

$\max=10$, $\sum=32$. Binary search converges to $\textit{maxSum}=18$: one
valid split is $[7,2,5]$ (sum $14$) and $[10,8]$ (sum $18$), giving largest
sum $18$; checking $17$ shows it forces a third part
($[7,2,5]=14$, then $[10]=10$ would leave $[8]$ separately, or
$[7,2,5,10]=24>17$ forces an earlier break, ultimately requiring $3$
parts), confirming $18$ is minimal.

\subsection*{C++ implementation}

\begin{lstlisting}
#include <bits/stdc++.h>
using namespace std;

int partsNeeded(vector<int>& nums, int maxSum) {
    int parts = 1;
    long long currentSum = 0;

    for (int x : nums) {
        if (currentSum + x > maxSum) {
            parts++;
            currentSum = 0;
        }
        currentSum += x;
    }
    return parts;
}

int splitArray(vector<int>& nums, int m) {
    int lo = *max_element(nums.begin(), nums.end());
    int hi = accumulate(nums.begin(), nums.end(), 0);

    while (lo <= hi) {
        int mid = lo + (hi - lo) / 2;
        if (partsNeeded(nums, mid) <= m) {
            hi = mid - 1;
        } else {
            lo = mid + 1;
        }
    }
    return lo;
}

int main() {
    vector<int> nums = {7, 2, 5, 10, 8};
    cout << "Minimized largest sum: "
         << splitArray(nums, 2) << endl;
    return 0;
}
\end{lstlisting}

\begin{complexitybox}
\textbf{Time Complexity.}
$O(n \log(\sum(\textit{nums})))$, identical to
Section~\ref{sec:bs-book-allocation}.
\textbf{Space Complexity.} $O(1)$ auxiliary space.
\end{complexitybox}

\begin{takeawaybox}
The fastest way to solve this problem is to recognize it as Book
Allocation in disguise. Building a small catalogue of templates -- rather
than a catalogue of individual problems -- is what lets binary search on
the answer scale to dozens of seemingly distinct sheet problems with only
a handful of underlying ideas.
\end{takeawaybox}
